\chapter{Methods}
\label{cap:methods}

% Each phase of the work adds its subsection here as it is implemented, so that
% the report grows at the same pace as the code.

\section{Pipeline overview}
\label{sec:pipeline}

\todo{One figure and one paragraph: audio, features, embedding, cosine
similarity, kNN graph, evaluation. Say that the embedding box is deliberately
isolated, so that replacing MFCCs with a trained network changes one
configuration file and nothing else.}

\section{Experiment framework}
\label{sec:framework}

One experiment is one YAML file. A configuration file inherits from a global
default file, a dataset file and an extractor file through a \texttt{defaults}
list, and the resulting document is hashed; that hash names both the cache
directory of every expensive artefact and the directory of the run. Datasets,
extractors, graph builders, relevance functions and evaluators are registered
under a name, so adding a model means adding one class and one configuration
file rather than editing the pipeline.

Every execution writes a directory holding the resolved configuration, the
metrics, the figures, the log, and a provenance record with the git commit,
whether the working tree was clean, the exact version of every dependency, the
seed, the platform and the time each stage took. A number in this report can
therefore be traced back to the code and the configuration that produced it.

\section{Audio preprocessing}
\label{sec:audio}

\todo{Decoding, resampling to 22\,050\,Hz, conversion to mono, and fixing the
length to 661\,500 samples so that every spectrogram has the same number of
frames. Log-mel spectrograms: 2048-point FFT, hop 512, 128 mel bands.}

\section{Feature extraction}
\label{sec:features}

\todo{The MFCC baseline: 20 coefficients plus deltas, mean and standard
deviation pooled over time, 80 dimensions. Explain why $c_0$ is discarded: it is
the frame energy, that is, the mastering loudness of the MP3 rather than the
timbre, and it would dominate the cosine similarity. Describe the random control
and say that it is reported next to every model.}

\section{Classical techniques}
\label{sec:classical-methods}

\todo{Describe the classical block (T3), evaluated on exactly the same items as
every other representation: \texttt{acoustic}, the broader descriptor set of
Tzanetakis and Cook (chroma, spectral contrast, centroid, bandwidth, rolloff,
ZCR, RMS, tempo, on top of MFCC, about 170 dimensions); \texttt{gauss\_mfcc}, a
single Gaussian per song over the MFCC frames (mean plus $20\times20$
covariance), compared with symmetric KL divergence; and the classical
classifiers (kNN, logistic regression, SVM-RBF, random forest) trained on
\texttt{acoustic}, used both as genre classifiers and, through their posterior
probability vector \texttt{posterior\_<clf>}, as an 8-dimensional
representation for graph construction. State the dimension control: every
representation is also evaluated after PCA to 64 dimensions, fitted on
training, except representations already smaller than 64.}

\section{Embeddings}
\label{sec:embeddings}

\todo{Standardisation fitted on the training split only, optional PCA, and L2
normalisation so that cosine similarity is an inner product. Explain why the
scaler is fitted on MTT itself rather than transferred from FMA.}

\section{Siamese networks}
\label{sec:siamese-models}

\todo{Phase R8/R9. This is the principal focus of the work (T2). Describe the
three siamese inputs that answer SQ5: \texttt{siamese\_mel} (log-mel
spectrogram, common CNN backbone), \texttt{siamese\_emb} (a frozen audio
embedding, the headline layer of MERT, through an MLP head) and
\texttt{siamese\_desc} (the acoustic descriptors of Section
\ref{sec:classical-methods}, through an MLP head), all trained with the
contrastive loss \cite{Hadsell2006} and producing 128-dimensional embeddings.
Add \texttt{triplet\_mel}, the same spectrogram input trained with the
semi-hard triplet loss \cite{Schroff2015}, to compare pairs against triplets.
State the sampling rule and its reason: a positive pair is two clips of the
same genre by different artists, because pairing clips of one artist teaches
the network to recognise the artist and the album production with it. State
the controls that make the siamese-input comparison valid: an equalised
training budget and hyper-parameter grid across \texttt{siamese\_*},
\texttt{triplet\_mel} and \texttt{cnn}; the same backbone for every
spectrogram-input model; comparable head capacity, with parameter counts
reported; and 3 seeds per trained model. State the headline result this
comparison targets: separating the gain a siamese network adds over its raw
input (\texttt{siamese\_emb} against \texttt{mert}, \texttt{siamese\_desc}
against \texttt{acoustic}) from the gain that comes from the input itself.}

\section{The CNN classifier and MERT}
\label{sec:trained-models}

\todo{Phase R8/R9. The convolutional backbone shared with \texttt{siamese\_mel},
trained here as a genre classifier with cross-entropy, with the embedding taken
from the penultimate layer. Then MERT \cite{Li2023}: 24\,kHz, 13 layers saved,
time-averaged per layer and per window, with the headline layer chosen on the
validation split only through a linear probe.}

\section{Similarity graph}
\label{sec:graph}

\todo{Brute-force cosine kNN in bounded-memory row blocks. Be explicit that
self-exclusion is done by row index rather than by comparing similarity values:
with duplicated vectors a value-based rule silently removes the wrong
neighbour. Describe the near-duplicate detection and the variants: mutual kNN,
the $\varepsilon$ threshold, and weighted against unweighted edges. The
canonical crisp graph is directed kNN by cosine, $k=10$; it is symmetrised by
union for community and clustering analyses.}

\section{Similarity measures}
\label{sec:similarity-measures}

\todo{Describe the measure registry (T3): cosine as the canonical measure,
Euclidean and Manhattan after standardisation, Pearson correlation,
Mahalanobis with Ledoit-Wolf shrinkage for \texttt{mfcc} and \texttt{acoustic},
and symmetric KL divergence between Gaussians for \texttt{gauss\_mfcc}. Note
that the distance a siamese network learns is cosine in its own embedding
space, so it coincides with the canonical measure.}

\section{Fuzzy block}
\label{sec:fuzzy-methods}

\todo{Phase R11 (T4), applied on top of any representation, so that the gain
from modelling gradual relationships can be related to the representation
itself (SQ6). Describe the four techniques:
\begin{itemize}
  \item \textbf{F1, fuzzy similarity graph.} Local membership
        $\mu_{ij} = \exp(-(d_{ij} - \rho_i)/\sigma_i)$ over the $k$ nearest
        neighbours, with $\rho_i$ the distance to the nearest neighbour and
        $\sigma_i$ calibrated so that $\sum \mu = \log_2 k$; symmetrisation
        with t-conorms (maximum, probabilistic sum, {\L}ukasiewicz); and
        $\alpha$-cuts to recover a crisp graph.
  \item \textbf{F2, fuzzy inference (if confirmed by the supervisor).} A small,
        interpretable rule base over facet similarities (timbre, harmony,
        rhythm, energy), with triangular or trapezoidal membership functions
        and Mamdani or zero-order Takagi-Sugeno inference; parameters fitted on
        the FMA validation split, before touching test.
  \item \textbf{F3, fuzzy genre membership.} Fuzzy kNN (Keller et al. 1985)
        over the graph, and fuzzy c-means (Bezdek 1981) over the embeddings,
        for overlapping genre communities.
  \item \textbf{F4, model-versus-human ambiguity.} A cheap analysis on top of
        the above: whether the model is uncertain where humans are uncertain,
        on the MTT triplets.
\end{itemize}
Implemented in NumPy, with no new dependency unless justified. State which of
F1--F4 are core and which depend on the supervisor's confirmation (open
decision in the project plan).}
